\scanpage{78}
C'est un système inductif de \deleted{\unsure{espaces}} proespaces
sur l'ens ordonné opposé de celui des drapeaux de $I$
(pour l'inclusion), dont le $\uline{\lim}$ est \unsure{égalt} $\approx X$.

Soit $Z$ \deleted{$X$} une partie (\inserted{loc\textsuperscript{t}}fermée)
de $X$, et posons, pour \deleted{$\unsure{\forall}$} un $i$ tel que
$X_i^*\subset Z$
\[
(55)\qquad V_i^Z=V_i\cap Z\qquad
\text{provoisinage tubulaire de }X_i^*\text{ dans }Z
\]
\deleted{de}
On a
\[
(56)\qquad X_i^*\hookrightarrow V_i^Z\hookrightarrow V_i
\qquad\text{équivalences d'homotopie}
\]
Plus généralement, on pose (si $X_{i_n}\subset Z$)
\[
V_{i_0-i_n}^Z=V_{i_0-i_n}\cap Z
=V_{i_0}^Z\cap\ldots\cap V_{i_n}^Z
\]
C'est aussi le $V_{i_0\ldots i_n}$ relatif à la \deleted{fil}
stratification de $Z$ induite par les
$Z_i\overset{\mathrm{déf}}{=}X_i\cap Z$, pour le système d'indices
$i_0\ldots i_n$. On a
\[
(57)\qquad V_{i_0\ldots i_n}^Z\hookrightarrow V_{i_0\ldots i_n}
\qquad\text{équiv. d'homotopie}
\]
Soit $J$ une partie de $I$, posons
\[
(58)\qquad X_J^*=\bigcup_{i\in J}X_i^*,\qquad X_J=\bigcup_{i\in J}X_i
\]
\mspage{77}
